\section{Open questions and next steps}

The five questions below survive this replication and each has a concrete,
free-compute experiment attached.

\begin{enumerate}
  \item \textbf{Does implementing the Bogda\l{} 2013 Gillespie apoptosis gate
        over a $\geq 100$-cell ensemble recover LUCID's Fig.~6 apoptosis
        percentages to within $\sim 10\%$ per dose bin?}
        This is the single largest lever to move from PARTIAL to REPLICATED.
        \emph{Next steps.} Port the Bogda\l{} 2013 apoptosis reactions from
        the PLOS Comp.\ Biol.\ S1 Text into a Gillespie / $\tau$-leap engine
        wrapped around the existing SciPy ODE substrate; run $\geq 100$
        independent cells per dose per Hill threshold $M$; compare bar
        heights to LUCID Fig.~6a ($M=0.5$~Gy) and Fig.~6b ($M=0.14$~Gy).
        Compute: single CPU, hours to overnight.

  \item \textbf{How does the p53 network behave under combined-modality
        insult (IR + nutlin-3a or IR + topo-II poison)?}
        The paper models radiation alone; combined-modality is the current
        translational-oncology frontier and may drive the system into
        regimes not covered by the Hat/Bogda\l{} parameter fits.
        \emph{Next steps.} Add a nutlin-3a term as a Mdm2--p53 binding-affinity
        multiplier (Vassilev 2004 $K_d$); sweep $[\mathrm{nutlin}] \times
        [\mathrm{dose}]$ on a $5\times5$ grid; report 72-h apoptosis fraction
        and time-to-peak-p53 as a heatmap; cross-check any published in-vitro
        combination-index data (MCF-7 or U2OS panels).

  \item \textbf{How sensitively do the predictions depend on the ATM Hill
        threshold $M$, and can $M$ be fit from single-cell live-imaging data
        rather than chosen between two literature values?}
        The replication shows apoptotic propensity is systematically higher
        at $M=0.14$~Gy than $M=0.5$~Gy, matching the \emph{sign} of LUCID
        Fig.~6a vs~6b, but the magnitude of that shift is unquantified.
        \emph{Next steps.} Fit $M$ against published p53-Venus live-imaging
        traces (Lahav lab; Purvis 2012; Batchelor 2011); run profile-likelihood
        or MCMC identifiability on the ATM Hill subset; propagate posterior
        uncertainty forward to the apoptosis prediction.

  \item \textbf{Does replacing the analytical DSB square pulse with a
        track-structure Monte Carlo change apoptotic fraction for high-LET
        beams (protons, $\alpha$, carbon)?}
        The current implementation is silent on LET because it inherits
        Hat 2016's analytical DSB input; LUCID's stated motivation is exactly
        to couple track-structure physics to the p53 network.
        \emph{Next steps.} Install TOPAS-nBio (free academic license) or use
        published open DSB-complexity distributions per LET; drive the ODE's
        DSB compartment with the resulting time-series; compare apoptosis
        fractions at matched dose but different LET; check whether the model
        predicts the known RBE elevation for high-LET beams.

  \item \textbf{Can a lightweight ML surrogate (GP or small MLP) trained on
        the LUCID ensemble predict per-cell fate from the first 4 h of a
        simulated p53 trace, and transfer to real Lahav-lab p53-Venus traces
        without re-training?}
        Live-cell p53 imaging produces rich timelapse data the paper does not
        exploit. A trained surrogate would let the mechanistic model act as a
        prior for cell-fate classification and sharpen the question of which
        single-cell features carry the fate signal.
        \emph{Next steps.} Generate $\sim 10^4$ simulated traces across a
        dose $\times M \times$ noise-level grid; label each with its 72-h fate
        from the stochastic gate; train a classifier on the first-4h window;
        validate on held-out simulations; apply to Lahav-lab p53-Venus traces
        (or DREAM Challenge data) and report classification accuracy and
        calibration.
\end{enumerate}
